% ============================================================
%  SECCION 2 — EXPLICACIONES TECNICAS DE HIPERPARAMETROS
% ============================================================
\section{Explicaciones Tecnicas de Hiperparametros}

\begin{tcolorbox}[colback=ujatblue!10, colframe=ujatblue,
  title={\bfseries Proposito de esta seccion}, coltitle=white]
Cada decision de diseno del experimento tiene una justificacion tecnica.
Esta seccion explica el \emph{por que} de cada hiperparametro en NB01 (1D)
y NB02 (2D) para que pueda responder con precision al comite.
\end{tcolorbox}

\vspace{8pt}

% ==============================================================
\subsection{Arquitectura comun SIREN}
% ==============================================================

\subsubsection*{Activacion sinusoidal: $\phi(z) = \sin(\omega_0 \cdot z)$}

\begin{tcolorbox}[colback=gray!5, colframe=ujatgray, breakable]
La ecuacion de Helmholtz requiere derivadas de segundo orden $\nabla^2 E$.
Para que la perdida de fisica sea util, la red debe ser diferenciable al menos
dos veces con derivadas no degeneradas.

\begin{tabular}{@{}lll@{}}
\toprule
\textbf{Activacion} & \textbf{Segunda derivada} & \textbf{Apta para Helmholtz?} \\
\midrule
ReLU & $= 0$ (casi en todos lados) & \textbf{No} \\
tanh & $\to 0$ para $|x| \gg 0$ & Muy limitada \\
$\sin(\omega_0 z)$ & $= -\omega_0^2 \sin(\omega_0 z)$ & \textbf{Si, siempre} \\
\bottomrule
\end{tabular}

\vspace{4pt}
Con SIREN, la $n$-esima derivada es $\pm \omega_0^n \sin(\omega_0 z + n\pi/2)$:
siempre una sinusoide, nunca cero, nunca saturada. Esto hace la perdida
$\mathcal{L}_\text{fisica} = \|\nabla^2 E + k^2 E\|^2$ bien condicionada.
\end{tcolorbox}

\subsubsection*{Inicializacion especial de pesos (Sitzmann et al., 2020)}

\begin{tcolorbox}[colback=gray!5, colframe=ujatgray, breakable]
Si se inicializa SIREN con Xavier o He estandar, las activaciones sinusoidales
se saturan o colapsan en las primeras epocas. La solucion es calibrar la varianza
de los pesos por capa:

\[
  W_0 \sim U\!\left(-\frac{1}{n}, +\frac{1}{n}\right)  \quad \text{(capa de entrada)}
\]
\[
  W_i \sim U\!\left(-\frac{\sqrt{6/n}}{\omega_0}, +\frac{\sqrt{6/n}}{\omega_0}\right)
  \quad \text{(capas ocultas)}
\]

donde $n$ es el fan-in de la capa. El factor $1/\omega_0$ compensa la amplificacion
de $\omega_0$ en la activacion: garantiza que $\sin(\omega_0 \cdot W_i x)$ explore
el rango $[-1, +1]$ uniformemente desde el inicio.

\textbf{Consecuencia practica:} sin esta inicializacion, la red converge a $E = 0$
o $E = 1$ constantes. Con ella, la red aprende soluciones oscilatorias desde
la primera epoca.
\end{tcolorbox}

\vspace{10pt}

% ==============================================================
\subsection{NB01: Helmholtz 1D --- Hiperparametros}
% ==============================================================

\begin{tcolorbox}[colback=slidelightblue!30, colframe=ujatblue,
  title={\bfseries Configuracion NB01}, coltitle=white, breakable]
\begin{tabular}{@{}lll@{}}
\toprule
\textbf{Hiperparametro} & \textbf{Valor} & \textbf{Justificacion} \\
\midrule
Arquitectura & SIREN 5$\times$64 & Ver abajo \\
Parametros totales & 16,833 & Suficiente para 1D sinusoidal \\
$\omega_0$ & 1.0 & Ver abajo \\
$\lambda_\text{fis}$ & 1.0 & 1D: residuo unico, sin overflow \\
$N_\text{colloc}$ & 2,000 (uniforme) & 1D: cobertura uniforme es suficiente \\
$N_\text{boundary}$ & 5 puntos & $x \in \{0, 0.25, 0.5, 0.75, 1\}$ \\
Paciencia Adam & umbral $\mathcal{L} < 10^{-4}$ & Criterio de parada por tolerancia \\
L-BFGS max\_iter & 500 & Problema 1D mas simple \\
L-BFGS history & 50 & Menor memoria: curvatura local suficiente \\
SEED & 42 & Reproducibilidad global \\
\bottomrule
\end{tabular}
\end{tcolorbox}

\vspace{6pt}

\subsubsection*{?`Por que arquitectura 5$\times$64 y no mas grande?}
\begin{tcolorbox}[colback=gray!5, colframe=ujatgray, breakable]
En 1D la solucion es $E(x) = \cos(kx)$, una sinusoide simple. 5 capas de 64 neuronas
(16,833 parametros) son suficientes para representar esta funcion y su derivada segunda.
Aumentar la anchura a 128 no mejora el resultado (ya se alcanza $L^2 = 0.006\%$) pero
si duplica el tiempo de entrenamiento. Regla practica: usar la arquitectura minima que
cumple la meta de precision.
\end{tcolorbox}

\subsubsection*{?`Por que $\omega_0 = 1.0$?}
\begin{tcolorbox}[colback=gray!5, colframe=ujatgray, breakable]
El paper original de Sitzmann usa $\omega_0 = 30$, calibrado para imagenes fotograficas
donde las frecuencias tipicas son altas (ciclos por pixel). En este trabajo:
\begin{itemize}
  \item Dominio: $[0,1]^2$ adimensional
  \item Numero de onda: $k = 2\pi$
  \item Regla de calibracion: $\omega_0 \approx k/(2\pi) = 1$
\end{itemize}

Con $\omega_0 = 30$ y $k = 2\pi$, la activacion $\sin(30 \cdot Wx)$ crea gradientes
con amplitud de orden $30^2 = 900$ veces mayor de lo necesario. En la practica esto
produce explosión de gradientes en las primeras epocas de Adam y la perdida no baja
del nivel inicial. El ajuste $\omega_0 = 1.0$ resuelve este problema completamente.

\textbf{Nota:} Esta calibracion fue descubierta experimentalmente en NB01 y luego
documentada como regla general para dominios adimensionales.
\end{tcolorbox}

\subsubsection*{?`Por que 5 puntos de frontera (y no solo 2)?}
\begin{tcolorbox}[colback=gray!5, colframe=ujatgray, breakable]
La condicion de frontera es $E(0) = 1$ y $E(1) = 1$. Si se usan solo dos puntos
de frontera, la red puede colapsar a la solucion constante $E(x) = 1$, que satisface
las BCs pero no la ecuacion de Helmholtz en el interior. Al agregar tres puntos
intermedios conocidos ($x = 0.25, 0.5, 0.75$ con sus valores exactos de $\cos(kx)$),
se rompe esa degeneracion y la red aprende la oscilacion correcta.

Esta fue una leccion critica aprendida durante el desarrollo de NB01.
\end{tcolorbox}

\subsubsection*{?`Por que $\lambda_\text{fis} = 1.0$ en 1D?}
\begin{tcolorbox}[colback=gray!5, colframe=ujatgray, breakable]
En 1D la funcion de perdida es:
\[
  \mathcal{L} = \mathcal{L}_\text{BC} + 1.0 \cdot \mathcal{L}_\text{fisica}
\]
donde $\mathcal{L}_\text{fisica} = \text{MSE}(E'' + k^2 E)$ es un solo residuo escalar.
Con $\lambda = 1.0$ el balance es adecuado y los gradientes no desbordan float32.
El problema de overflow solo aparece en 2D (ver siguiente seccion).
\end{tcolorbox}

\vspace{10pt}

% ==============================================================
\subsection{NB02: Helmholtz 2D --- Hiperparametros}
% ==============================================================

\begin{tcolorbox}[colback=slidelightblue!30, colframe=ujatblue,
  title={\bfseries Configuracion NB02}, coltitle=white, breakable]
\begin{tabular}{@{}lll@{}}
\toprule
\textbf{Hiperparametro} & \textbf{Valor} & \textbf{Justificacion} \\
\midrule
Arquitectura & SIREN 5$\times$128 & Mayor capacidad para campo complejo \\
Parametros totales & 66,690 & 4$\times$ NB01: 2 salidas, variacion 2D \\
$\omega_0$ & 1.0 & Misma regla que NB01 \\
$\lambda_\text{fis}$ & \textbf{0.1} & \textbf{Critico: 2 residuos duplican peso fisico} \\
$N_\text{colloc}$ & 3,000 LHS & LHS previene clustering en 2D \\
$N_\text{boundary}$ & 1,200 total & 300 por borde $\times$ 4 bordes \\
Paciencia Adam & 800 epocas & Problema 2D mas dificil \\
L-BFGS max\_iter & 1,000 & Mayor exigencia de convergencia \\
L-BFGS history & 100 & Mayor memoria para curvatura 2D \\
SEED & 42 & Reproducibilidad global \\
\bottomrule
\end{tabular}
\end{tcolorbox}

\vspace{6pt}

\subsubsection*{El cambio critico: $\lambda_\text{fis} = 0.1$ (no $1.0$)}

\begin{tcolorbox}[colback=slidewarning!60, colframe=ujatorange,
  title={\bfseries Explicacion de la calibracion de $\lambda$}, coltitle=white, breakable]
\textbf{?`Por que no usar $\lambda = 1.0$ como en 1D?}

En 2D la funcion de perdida tiene \textbf{dos residuos} (parte real e imaginaria):
\[
  \mathcal{L}_\text{fisica} = \text{MSE}(R_\text{real}^2) + \text{MSE}(R_\text{imag}^2)
\]
donde $R_\text{real} = \nabla^2 E_r + k^2 E_r$ y $R_\text{imag} = \nabla^2 E_i + k^2 E_i$.

Esto \textbf{duplica el peso efectivo} de la fisica:
\[
  \mathcal{L} = \mathcal{L}_\text{BC} + \lambda \cdot (\mathcal{L}_\text{real} + \mathcal{L}_\text{imag})
\]
Con $\lambda = 1.0$: el peso efectivo de la fisica es $\approx 2.0$. En float32 (precision
de 32 bits), los gradientes de la perdida fisica alcanzan valores del orden $10^{38}$
y producen NaN $\to$ CUDA crash.

\textbf{Ablacion verificada} (\code{results/ablation\_lambda.json}):
\begin{center}
\begin{tabular}{@{}lccc@{}}
\toprule
$\lambda$ & Error $L^2$ & Estado & Observacion \\
\midrule
0.01 & 0.163\% & OK & Sub-ponderacion de la fisica \\
\textbf{0.1} & \textbf{0.171\%} & \textbf{Optimo} & Balance fisica/frontera \\
1.0 & --- & CRASH & Overflow float32 en CUDA \\
\bottomrule
\end{tabular}
\end{center}

Con $\lambda = 0.01$ la fisica queda sub-ponderada: la red ajusta bien las BCs
pero no maximiza la satisfaccion de Helmholtz. Con $\lambda = 0.1$ se logra el
balance optimo.
\end{tcolorbox}

\subsubsection*{?`Por que LHS en lugar de muestreo uniforme?}
\begin{tcolorbox}[colback=gray!5, colframe=ujatgray, breakable]
Latin Hypercube Sampling (LHS) estratifica el espacio de entrada:
\begin{itemize}
  \item Divide cada dimension en $N = 3000$ intervalos iguales
  \item Garantiza exactamente un punto en cada intervalo por dimension
  \item El resultado: los puntos de colocacion cubren el dominio sin clusters ni
        zonas vacias, independientemente de la semilla aleatoria
\end{itemize}

El muestreo uniforme aleatorio en 2D puede concentrar puntos en algunas zonas y
dejar otras sin cobertura. En una PINN esto se traduce en mayor error en las zonas
poco muestreadas. LHS elimina este problema de forma garantizada.

\textbf{Nota:} en 1D el muestreo uniforme es suficiente. LHS solo agrega valor
a partir de 2 dimensiones.
\end{tcolorbox}

\subsubsection*{?`Por que arquitectura 5$\times$128 (y no 5$\times$64)?}
\begin{tcolorbox}[colback=gray!5, colframe=ujatgray, breakable]
En 2D hay tres razones para ampliar la arquitectura:
\begin{enumerate}
  \item \textbf{Campo complejo}: la red predice dos salidas simultaneas ($E_r$ y $E_i$).
        Esto requiere mayor capacidad para representar las correlaciones entre ambas partes.
  \item \textbf{Variacion 2D}: la solucion varia en dos dimensiones; la red necesita
        mas neuronas para capturar la frecuencia espacial en ambas direcciones.
  \item \textbf{Numero de colocation points}: 3,000 puntos en 2D vs 2,000 en 1D.
        La red debe poder interpolar correctamente entre mas restricciones.
\end{enumerate}
Con 5$\times$64 en 2D el error L2 sube a $\approx$0.8\% (estimado). Con 5$\times$128
se alcanza 0.171\%. Aumentar a 5$\times$256 no mejora significativamente pero
cuadruplica el tiempo de entrenamiento.
\end{tcolorbox}

\subsubsection*{?`Por que Adam primero y luego L-BFGS?}
\begin{tcolorbox}[colback=gray!5, colframe=ujatgray, breakable]
\textbf{Adam} (Adaptive Moment Estimation):
\begin{itemize}
  \item Algoritmo estocástico de primer orden
  \item Estima el gradiente con mini-batches $\to$ navega paisajes de perdida no convexos
  \item Converge hacia la region del minimo global en 10,000--15,000 epocas
  \item Limitacion: precision limitada, no usa curvatura (Hessiano)
\end{itemize}

\textbf{L-BFGS} (Limited-memory Broyden-Fletcher-Goldfarb-Shanno):
\begin{itemize}
  \item Metodo cuasi-Newton de segundo orden
  \item Aproxima el Hessiano con un historial de gradientes (history\_size=100)
  \item Convergencia cuadratica cerca del minimo: cada iteracion reduce el error
        por un orden de magnitud
  \item Requiere un punto inicial ya cercano al minimo (si se usa desde cero diverge)
  \item Usa busqueda de linea strong Wolfe: garantiza descenso en cada iteracion
\end{itemize}

\textbf{Sinergia}: Adam lleva los pesos a la ``cuenca de atraccion'' del minimo,
L-BFGS los lleva al minimo preciso. Usar solo Adam da $L^2 \approx 1\%$; la fase
L-BFGS lo reduce a $L^2 = 0.171\%$.

\textbf{Leccion critica}: nunca usar \emph{learning rate scheduler} antes de L-BFGS.
Si Adam usa scheduler (decaimiento de LR), los pesos llegan a L-BFGS en un estado
``congelado'' donde L-BFGS solo logra 7/1000 iteraciones. Sin scheduler, L-BFGS
completa 1,035 iteraciones y converge a la precision esperada.
\end{tcolorbox}

\subsubsection*{?`Por que paciencia=800 epocas en NB02 y no un umbral fijo?}
\begin{tcolorbox}[colback=gray!5, colframe=ujatgray, breakable]
En NB01 se usa un umbral de tolerancia ($\mathcal{L} < 10^{-4}$) porque la solucion
1D converge rapidamente y de forma predecible. En NB02 la perdida en 2D fluctua mas
durante Adam; un umbral fijo puede terminar prematuramente el entrenamiento en una
meseta temporal. La paciencia de 800 epocas garantiza que Adam explora el espacio
suficientemente antes de pasar a L-BFGS, independientemente del nivel de perdida.
\end{tcolorbox}

\vspace{10pt}

% ==============================================================
\subsection{Validacion Multi-Semilla}
% ==============================================================

\begin{tcolorbox}[colback=slidelightgreen!40, colframe=ujatgreen,
  title={\bfseries Por que tres semillas?}, coltitle=white, breakable]
Las redes neuronales son sensibles a la inicializacion aleatoria. Un resultado con
una sola semilla puede ser un ``accidente afortunado''. La validacion multi-semilla
cuantifica si el metodo es robusto:

\begin{center}
\begin{tabular}{@{}ccc@{}}
\toprule
\textbf{Semilla} & \textbf{Error $L^2$ (\%)} & \textbf{Epocas} \\
\midrule
42  & $\approx$0.171 & $\approx$15,000 \\
123 & $\approx$0.148 & $\approx$14,200 \\
777 & $\approx$0.146 & $\approx$14,500 \\
\midrule
\textbf{Promedio} & \textbf{$0.155 \pm 0.020$} & \\
\bottomrule
\end{tabular}
\end{center}

La desviacion estandar del 0.020\% es dos ordenes de magnitud menor que el umbral
del 5\%. Esto confirma que el metodo converge de manera consistente, no dependiente
de la inicializacion particular.

\textbf{Fuente verificada}: \code{results/multiseed\_results.json}
\end{tcolorbox}

\vspace{10pt}

% ==============================================================
\subsection{Resumen de Decisiones de Diseno}
% ==============================================================

\begin{table}[h!]
\centering
\small
\begin{tabular}{llll}
\toprule
\textbf{Decision} & \textbf{NB01} & \textbf{NB02} & \textbf{Razon del cambio} \\
\midrule
Anchura de la red & 64 neuronas & 128 neuronas & Campo complejo + 2D \\
$\lambda_\text{fis}$ & 1.0 & \textbf{0.1} & 2 residuos $\to$ overflow float32 \\
Muestreo & Uniforme & \textbf{LHS} & LHS solo vale en $\geq$ 2D \\
$N_\text{colloc}$ & 2,000 & 3,000 & Mayor densidad en 2D \\
$N_\text{boundary}$ & 5 pts & 1,200 pts & 4 bordes completos \\
Paciencia Adam & Umbral L & 800 epocas & Fluctuaciones en 2D \\
L-BFGS max\_iter & 500 & 1,000 & Convergencia mas exigente \\
\bottomrule
\end{tabular}
\caption{Diferencias entre NB01 y NB02. El cambio mas critico:
$\lambda_\text{fis}$ de 1.0 a 0.1 (evita overflow float32 en 2D).}
\end{table}
